% Section B.9, from lectures/L09/appendix/c_exercises.md: every part that is not code. The outcomes
% quoted were measured on a reference build of the whole stack made to the chapters'
% specifications: the host suites over the mocked register file, and the avr-gcc 7.3.0 image and
% its disassembly.

\section{Chapter 9: The Real Transport, and Both Toolchains}
\label{sol:c9}
The code in this chapter is \code{AvrSpi} (Exercise~\ref{c9:ex:1} a), checked by the provided
suite \file{fw/test/transport/avr_spi_test.cpp} over the mocked register file, and the bring-up
\code{main} and the logging (Exercise~\ref{c9:ex:2} a and Exercise~\ref{c9:ex:3} a), which the
bench checks. Everything else is answered here.

% ------------------------------------------------------------------------------------------
\solution{c9:ex:1}{\code{AvrSpi}}
\ans{b} \code{SPDR} holds the incoming byte only once all eight bits have been shifted, eight
\code{SCK} periods after the write, which at f\textsubscript{osc}/16 is 128 CPU cycles, and
\code{SPIF} is the only thing that says that moment has come: read earlier and you get a stale
byte, write the next byte earlier and you collide with the transfer still in progress.

The spin is cheaper than it looks. Compiled, the whole of \code{transfer()} is six instructions,
and the wait is a four-cycle loop that goes round about thirty times per byte:

\begin{console}
out  0x2e, r22      ; SPDR = byte: the transfer starts
in   r0, 0x2d       ; read SPSR
sbrs r0, 7          ; SPIF set? skip the jump
rjmp .-6            ; not yet: read SPSR again
in   r24, 0x2e      ; return SPDR
ret
\end{console}

\ans{c} If \code{SS} (PB2) is an input and something drives it low while \code{MSTR} is set, the
SPI peripheral concludes that another master has taken the bus and gets out of its way:
\textbf{it clears \code{MSTR} in \code{SPCR}}, demoting itself to a slave, and sets \code{SPIF},
raising the SPI interrupt if one is enabled. The ATmega328P has no separate mode-fault flag;
\code{MSTR} reading back as 0 is the only record of what happened.

For \code{AvrSpi} that is two failures in a row. The \code{transfer()} in progress sees \code{SPIF}
and returns whatever is in \code{SPDR}. Every \code{transfer()} after it writes \code{SPDR} as a
slave, where nothing is driving \code{SCK}, so \code{SPIF} never sets again and the poll spins for
ever: the driver hangs in the first transaction after the glitch, and nothing sets \code{MSTR} back.
Making PB2 an output removes the possibility, and costs nothing, because the transport drives that
pin as the chip select anyway.

\ans{d} It returns the byte the \emph{previous} transfer received, still sitting in \code{SPDR}'s
receive buffer, or whatever was there after reset if there was no previous transfer. The new byte
does not exist yet: it arrives eight \code{SCK} periods later, 8\,\textmu s or 128 CPU cycles after
the write, and the read happens one cycle after it. Every reply is therefore shifted by one byte,
so a \code{readReg} assembles the command byte's reply and three data bytes instead of four.

On the part it is worse than a shift. The next \code{transfer()} writes \code{SPDR} a few cycles
later, while the byte before it is still shifting, which the datasheet calls a write collision:
\code{WCOL} is set and the write is \emph{ignored}, so most of the bytes of a transaction never
leave the master at all, and the five-byte transaction is garbled in both directions.

The mock makes the same point more bluntly. It completes a transfer only when the driver reads
\code{SPSR}, so with the poll gone no transfer ever completes and every read of \code{SPDR} returns
the value left from reset. Three cases go red, and only these three:

\begin{console}
Test case AvrSpi.TransferExchangesByte failed: EXPECT_EQ(rx, miso) failed: 0 != 205
Test case AvrSpi.TransferIsFullDuplexInOrder failed: EXPECT_EQ(spi.transfer(mosi[i]), miso[i])
  failed: 0 != 170
Test case AvrSpi.DrivesTheDriverEndToEnd failed: EXPECT_EQ(uart.status(), data) failed:
  0 != 305419896
19 out of 22 test cases succeeded!
\end{console}

\noindent Every \code{Uart} case stays green, because those run over the scripted
\code{transport::Stub}, not over \code{AvrSpi}: the bug is below the seam, and so is the only suite
that can see it.

% ------------------------------------------------------------------------------------------
\solution{c9:ex:2}{The freestanding bring-up \code{main}}
\ans{b} Not the provided demo itself: \file{source/main.cpp} was written in the AVR-portable
subset, and it compiles and links unchanged with the \file{fw/avr} flags, into 1256 bytes of flash
and 34 of RAM. It needs none of the idioms a freestanding build removes: its objects are automatic,
with fixed arrays inside; the driver reports failure with a \code{bool}; and it prints nothing.

It is the idioms of ordinary host C++ that fail, and each has a replacement on the AVR:

\begin{booktable}{@{}p{10.6em}p{11.4em}L@{}}
  \thead{Host idiom} & \thead{In the AVR build} & \thead{What replaces it} \\
  \midrule
  \code{#include <cstdint>}, \code{std::uint8_t} & Does not compile: avr-libc has no C++ library,
    so \code{fatal error: cstdint: No such file or directory} & \code{<stdint.h>} and the bare
    \code{uint8_t} \\
  \code{new}, \code{std::make_unique}, \code{std::vector} & No \code{operator new}, and no standard
    containers & Static and automatic storage, sized at compile time \\
  \code{throw}, \code{try} & Rejected under \code{-fno-exceptions} & Return values, as
    \code{write()} returning \code{false} \\
  \code{dynamic_cast}, \code{typeid} & Rejected under \code{-fno-rtti} & Virtual dispatch through
    an interface \\
  \code{std::cout}, \code{printf} to a console & No iostreams, and no console & Bytes written to the
    ATmega's own USART, as in Exercise~\ref{c9:ex:3} \\
\end{booktable}

\ans{c} With the reference firmware, the linker reports one symbol, three times:

\begin{console}
In function `app::EchoNode::~EchoNode()':
  undefined reference to `operator delete(void*, unsigned int)'
In function `driver::transport::AvrSpi::~AvrSpi()':
  undefined reference to `operator delete(void*, unsigned int)'
In function `driver::uart::Uart::~Uart()':
  undefined reference to `operator delete(void*, unsigned int)'
\end{console}

\noindent The construct is the \textbf{virtual destructor}. Every class that has one gets a second,
``deleting'' destructor, which runs the destructor and then calls \code{operator delete}, so that
\code{delete} through a base pointer frees the right object; the vtable points at it, so it is
emitted, and names \code{operator delete}, even though nothing in the program ever deletes anything.
From C++14 it calls the sized form, \code{operator delete(void*, size_t)}, and \code{size_t} is an
\code{unsigned int} on the AVR. That is why the linker names the three concrete classes: each
inherits a virtual destructor from its \code{Interface}.

The other symbols \file{env.cpp} defines turn up with other flags and other code, which is why it
defines them all. At \code{-O0} the same link also misses \code{__cxa_pure_virtual}, eight times:
the \code{Interface} classes' own vtables name it for every pure virtual, and only the optimizer
removes those vtables, because every \code{Interface} is constructed as the base of a concrete class
that overwrites its vtable pointer at once. Without \code{-fno-threadsafe-statics}, a
function-local \code{static} with a constructor, such as a \code{static driver::uart::Uart
uart\{spi\};} inside \code{main()}, adds \code{__cxa_guard_acquire} and \code{__cxa_guard_release}.
And under \code{-fno-sized-deallocation} the missing symbol is the unsized
\code{operator delete(void*)} instead.

The host build does not need the file because g++ links libstdc++, which defines every one of these
symbols, its static guards thread-safe. \file{env.cpp} lives in \file{avr/}, where the host Makefile
never looks, precisely so that its single-threaded guards can never replace those.

% ------------------------------------------------------------------------------------------
\solution{c9:ex:3}{Logging and timing}
\ans{b} Five bytes of eight bits is \textbf{40 \code{SCK} cycles}, which at 1\,MHz is \textbf{40\,\textmu
s} of clocking. The analyzer shows more: five bursts of eight clocks, each 8\,\textmu s long, with a
gap between consecutive bursts, and \code{SS} falling a little before the first and rising a little
after the last.

The gaps are software. \code{SCK} runs only while a byte is shifting, and between two bytes the CPU
has to notice \code{SPIF}, return from \code{transfer()}, go round \code{readReg}'s loop and fold the
byte into its 32-bit result, make the next virtual call through the transport interface, and write
\code{SPDR} again. Counted from the reference build's disassembly, that path is about 30 CPU cycles,
roughly 2\,\textmu s at 16\,MHz, so a register read takes about $5 \times 8 + 4 \times 2 = 48$\,\textmu
s between the first clock and the last, and about 50\,\textmu s from \code{SS} falling to \code{SS}
rising, once the \code{begin()} and \code{end()} calls on either side are counted. Your gaps depend on
your code and your compiler, but not by more than a microsecond or two; what does not change is that
the bus clocks for about four fifths of a transaction and waits for the CPU for the rest.

With the logging from part a) in the loop, the gaps \emph{between} transactions are a different
matter: a twenty-character line at 115200~baud takes 1.7\,ms to leave the USART, more than thirty
register reads.

\ans{c} The SPI is \textbf{synchronous}. The master's hardware makes \code{SCK} itself, by dividing
the CPU clock by the prescaler chosen in \code{SPCR} and \code{SPSR}, and sends it on its own wire;
the slave samples on the edges it receives. The code chooses only a ratio, f\textsubscript{osc}/16,
never a frequency, and the exact rate does not matter to either end so long as it stays within what
the slave and the wiring can take. No number in the transport depends on the clock, so it needs no
\code{F_CPU}.

The USART is \textbf{asynchronous}. There is no clock wire: the PC's end has a clock of its own and
agrees with the ATmega only on a number, the baud rate, which the ATmega has to produce from its own
clock by computing a divider, \code{UBRR0 = F_CPU / (16 * baud) - 1}, rounded to an integer. The code
can only compute that from the clock frequency, so \code{F_CPU} has to be the real one: with a wrong
\code{F_CPU} every character leaves at the wrong rate and arrives as garbage.

The rounding matters at the log's end of the line too. At 16\,MHz, 115200~baud needs a divider of
7.68, which rounds to 8 and gives 111111~baud, 3.5\,\% slow, at the edge of what a frame tolerates.
Setting \code{U2X0}, which halves the divisor's factor from 16 to 8, gives a divider of 16 and
117647~baud, 2.1\,\% fast; 38400~baud (divider 25) and 9600~baud (divider 103) are both within
0.2\,\%.

% ------------------------------------------------------------------------------------------
\solution{c9:ex:4}{Your own \code{uart_top} on the board}
\ans{a} The \textbf{Fitter} report's resource usage summary gives the size: the logic utilization
in ALMs, out of the 18\,480 of the DE0-CV's Cyclone~V (a 5CEBA4F23C7), the registers, and the pins.
This peripheral is a small fraction of the device, well under a few per cent; the exact figure
depends on your code and on the Quartus version, and is worth recording mainly so that you notice
if a later change makes it jump.

The \textbf{Timing Analyzer} gives the other number. Read the worst-case \textbf{setup slack} and
\textbf{hold slack} for the 50\,MHz clock across all the timing corners Quartus analyzes, slow and
fast silicon at the hot and the cold end of the temperature range; the multicorner summary gathers
the worst of them in one table. \textbf{Positive slack on every path, in every corner, means the
design met timing}; the restricted \code{Fmax} for the clock, which should come out comfortably
above 50\,MHz for a design this shallow, says the same thing another way. Check first that the
clock is actually constrained, a 20\,ns period on the 50\,MHz pin in the project's timing
constraints: an unconstrained clock is analyzed against nothing, and its slack means nothing. The
asynchronous inputs, \code{rx} and the three SPI inputs, are expected to show up as unconstrained,
because they enter through synchronizers and have no timing relationship to the clock to check.

\emph{The report} is how you know, and the design working is not, because the analysis covers the
slowest silicon, the lowest supply voltage and the highest temperature the part is specified for,
and the board on your desk is one sample, at room temperature, today. A design with negative slack
can work on the first try and fail warm, or on the next board; and nothing a working design does
tells you how much margin it has left.

\ans{c} The pin loopback proves the \textbf{physical path}, which the simulation never had in it:
that the FPGA is configured at all, that the wrapper's pin assignment puts \code{rx} and \code{tx} on
the header pins the wires go to, and that the adapter really is at 3.3\,V logic. It cannot tell you that the terminal's rate is
right: the adapter sends and receives at the same setting, so a loopback agrees with itself at any
rate, and the rate is first on trial when the peripheral talks to the terminal in \chapref{10}. No simulation has a pin, a voltage, a USB adapter
or a terminal in it.

\code{uart_top_tb} proves the \textbf{peripheral}, which the pin loopback never had in it: that a
byte written to \code{TX_DATA} over SPI leaves \code{uart_tx} correctly framed at the rate
\code{BAUD_DIV} sets, and comes back through \code{sync}, \code{uart_rx} and the RX FIFO to be read
from \code{RX_DATA}. With \code{rx} wired straight to \code{tx} in the wrapper, none of your logic is
in the path, so a character echoing in the terminal says nothing about it. Neither of the two says
anything about timing, which is what part a) is for.
